\section{Object Groups Approach} \label{sec:APPROACH}

We propose extending the Lively tools to allow editing multiple
objects at a time. Different objects' responsibilities overlap, making
it useful for them to share behavior. Objects that have multiple
responsibilities act in different \emph{roles} throughout the
system. Multiple-inheritance, traits, and layers may be used to
compose roles, but require either planning or refactoring into those
abstractions.

In our approach, developers can combine existing objects into
groups. Each group represents a specific responsibility, which objects
can be dynamically assigned to or removed from. This way developers
can develop concrete objects with immediate feedback, and modularize
them into groups on-demand to improve maintainability.

\subsection{Shared Properties and Shared Behavior}
\begin{itemize}
 \item snippets/evaluations are executed for every group member
 \item esp. addScript to add shared behavior
 \item this binds to every member of a group
 \item resultsets instead of single result when evaluating
 \item thus chages are propagating to every group member
\end{itemize}

\subsection{how to group objects}
\begin{itemize}
 \item select through magnifier
 \item select through scene-graph list
 \item through selection morph (WIP by Lauritz)
 \item edit()-function or ctrl-shift-e shortcut on a collection of morphs
\end{itemize}

\subsection{connection between tool and object}
\begin{itemize}
 \item magnifier to select and see selection (to connect tool to object)
 \item magnifier to browse morphic scene (hover over morph while selecting -- blue rectangle)
 \item tree widget instead of list - again: see and select (and select groups)
\end{itemize}

%%% Local Variables: 
%%% mode: latex
%%% TeX-master: "../paper"
%%% End: 
